\chapter{Working at Height Risk Assessment}
\label{chap:Working at Height Risk Assessment}

%---------------------------- SECTION ----------------------------
\section{Why this assessment matters in construction}

Falls from height remain the single largest cause of fatal injury in the UK construction industry, accounting for approximately one quarter of all construction fatalities in any given year. The Working at Height Regulations 2005 came into force following sustained HSE concern that the voluntary and prescriptive measures then in place were insufficient to reduce what was, and remains, a preventable category of death and serious injury. Despite the clarity of the regulatory framework, the comprehensiveness of HSE guidance, and the decades of industry awareness campaigns, operatives continue to fall from scaffolds, through fragile roofs, from ladders, and over unprotected edges. The persistence of this failure demands that every working at height risk assessment be conducted with meticulous care and that the comfortable assumptions that characterise inadequate assessments --- ``it's only a short duration job'', ``he's an experienced roofer'', ``the scaffold's been up for weeks without a problem'' --- be explicitly challenged.

The Working at Height Regulations 2005 define working at height as work in any place, including a place at or below ground level, from which a person could fall a distance likely to cause personal injury. This definition is broader than many practitioners assume: it includes work on flat roofs with unprotected edges, work on mezzanine floors without guardrails, and work in excavations with unsupported sides from which a person could topple. The duty to avoid work at height where reasonably practicable, to use collective protective measures in preference to individual protective measures, and to use personal protective equipment only where neither avoidance nor collective protection is practicable, reflects the hierarchy of controls applied specifically to the falls domain.

Falls of objects are frequently overshadowed by falls of persons in working at height assessments, but they represent a significant cause of injury and death in their own right. A falling tool, tile, or brick from a height of even three metres can cause a fatal head injury. The assessment of falling object risk must address the working area above the point where other trades or members of the public are present, the adequacy of toe-boards and netting on scaffold platforms, the management of materials and waste at height, and the provision of exclusion zones or covered walkways below areas of elevated work.

No working at height risk assessment is complete without a rescue plan. The Working at Height Regulations 2005 require that arrangements are made for the rescue of persons in the event of a fall, and this requirement is routinely neglected. A scaffold operative who falls while wearing a harness and is arrested by the lanyard before impact is in immediate danger from suspension trauma if not rescued within minutes. An operative working in a MEWP that suffers a mechanical failure is at risk of hyperthermia and suspension trauma if stranded at height. The rescue plan must be specific, pre-arranged, and tested; it must identify who will conduct the rescue, what equipment is available, and how quickly rescue can be initiated.

\barquote{``Falls from height are not accidents --- they are the predictable result of inadequate planning, inadequate equipment, and the tolerance of unsafe working practices that would never be accepted in any other context.''}

\example{Site Reality}{A roofing contractor is engaged to replace the felt covering on a flat commercial roof of approximately 400\,m$^2$. The roof has a central skylight bank comprising five rooflights, three of which are of uncertain structural integrity. The perimeter of the roof has a 300\,mm upstand parapet that provides no effective edge protection. The contractor's RAMS specify ``safety harness to be worn at all times'' as the primary fall protection measure, but the RAMS identifies no anchor points on the roof structure, specifies no rescue procedure, and does not address the fragile rooflights. Two operatives work for three days on the roof wearing harnesses clipped to each other's belt loops --- a practice that provides no fall arrest capability whatsoever. On the fourth morning, one operative steps onto a rooflight that fractures, falling 4.2 metres to the floor below. The absence of a suitable and sufficient working at height risk assessment --- with specific fragile roof protection, specific anchor points, and a specific rescue plan --- is identified as a primary causal factor.}

%---------------------------- SECTION ----------------------------
\section{Legal and professional framework}

The Working at Height Regulations 2005 are the primary legislative instrument for this hazard domain, applying to all work at height where there is a risk of a fall liable to cause personal injury. They impose a hierarchy of duty: first, avoid work at height where reasonably practicable; second, where work at height cannot be avoided, use work equipment or other measures to prevent a fall; third, where a fall cannot be prevented, use equipment or measures to minimise the distance and consequences of a fall. The Regulations are supported by PUWER 1998, which governs the suitability, maintenance, and inspection of work equipment used at height (including scaffolds, MEWPs, and ladders), and by LOLER 1998, which applies to MEWP thorough examinations and to any lifting operation conducted from height.

HSE guidance document INDG401 (Working at Height: A Brief Guide) provides accessible compliance guidance for those undertaking working at height risk assessments. BS EN 131 sets standards for portable ladders used in construction. The Construction Phase Plan required by CDM 2015 must specifically address the arrangements for managing working at height risks on the project, and the pre-construction information must identify any known fragile surfaces, existing edge protection, or structural limitations relevant to working at height.

\paragraph{Key frameworks relevant to this chapter include: Working at Height Regulations 2005; PUWER 1998; LOLER 1998; HSE INDG401; BS EN 131.}

\begin{itemize}
  \bitem{Working at Height Regulations 2005}{Imposes the hierarchy of duty to avoid, prevent, and minimise falls; requires that work at height is properly planned, appropriately supervised, and carried out safely; requires that the place of work at height is safe, that equipment is suitable and inspected, and that rescue arrangements are in place. Regulation 4 requires that no person engages in work at height unless competent to do so.}
  \bitem{PUWER 1998}{Requires that all work equipment used at height --- scaffolds, tower scaffolds, podium steps, MEWPs, ladders, and roof ladders --- is suitable for the task, maintained in an efficient state, inspected at appropriate intervals, and used only by trained and competent operatives.}
  \bitem{LOLER 1998}{Applies to MEWPs and to any lifting operation conducted from a working at height platform; requires thorough examination of MEWPs by a competent person at six-monthly intervals and planning of lifting operations by a competent person.}
  \bitem{BS EN 131 and scaffold standards}{BS EN 131 specifies the design, construction, and testing requirements for portable ladders; scaffold erection must comply with TG20:21 (NASC guidance for tube and fitting scaffold) or the manufacturer's guidance for system scaffold; scaffold must be inspected before first use, after any substantial alteration, and at a minimum of every seven days thereafter.}
\end{itemize}

%---------------------------- SECTION ----------------------------
\section{Conducting the assessment using the five-step method}

\subsection{Step 1 --- Identify hazards}
\paragraph{Working at height hazards in construction span a broad spectrum. Falls from height include: operative falls from unprotected scaffold edges; falls through fragile roof surfaces including rooflights, corrugated asbestos cement sheets, and aged metal decking; falls from unsecured or incorrectly angled ladders; falls from MEWPs due to overreaching, ground instability, or mechanical failure; and falls from leading edges of floor structures during construction. Falling objects include: tools and materials dropped from scaffold platforms; tiles, bricks, and masonry dislodged by working activity; and construction waste improperly disposed of from height. Scaffold failures include: collapse due to overloading, inadequate ties, or ground settlement; and failure of scaffold boards or platform components. The absence of a rescue plan following a fall-arrest event is itself an identifiable hazard.}

\subsection{Step 2 --- Decide who or what is affected}
\paragraph{Persons at risk from working at height hazards include: the operatives performing work at height; trades working beneath the elevated working area who are exposed to falling objects; scaffold erectors and dismantlers who are themselves at risk during the erection and dismantling phases; vehicle and plant operators whose cabs may be struck by falling objects; and members of the public beneath scaffold over public footpaths or within proximity of external scaffold. Vulnerable groups include: workers who have not received specific working at height training; young workers and apprentices with limited experience of height work; and workers with medical conditions (vertigo, cardiovascular conditions) that may affect their capacity to work safely at height.}

\subsection{Step 3 --- Evaluate likelihood and consequence}
\paragraph{Falls from height carry consequence scores at the maximum end of any risk matrix: fatality or life-changing injury is the realistic outcome of a fall of more than two metres. The inherent risk rating for unprotected working at height on a construction project must be assessed as Very High (typically 25 on a 5×5 matrix) before any controls are applied. The effectiveness of controls in reducing this inherent risk varies significantly: a fully-boarded, guardrailed scaffold with double handrails and toe-boards, regularly inspected, is a highly effective engineering control that can reduce the residual risk to Low; a harness and lanyard without confirmed anchor points, no rescue plan, and inadequate supervision reduces the inherent risk by a negligible amount and should not be scored as providing significant risk reduction.}

\subsection{Step 4 --- Select and implement controls}
\paragraph{Controls for working at height must be selected by applying the regulatory hierarchy: avoid the work at height where reasonably practicable (use ground-level assembly, pre-fabricated modules, long-reach equipment); prevent falls through collective protection (scaffold with full guardrail and toe-board, MEWP with integral protection, safety nets or airbags beneath the work area); minimise the distance and consequences of a fall through personal fall protection (harnesses, lanyards, and inertia reels connected to confirmed anchor points with a specific rescue plan). Every working at height activity must have a specific rescue plan attached to its risk assessment, naming the rescue team, the rescue equipment, and the procedure. Scaffold must be erected, altered, and dismantled only by competent persons and must be inspected at statutory intervals by a competent person with the inspection record retained on site.}

\subsection{Step 5 --- Record and review}
\paragraph{The working at height risk assessment must be reviewed before each day's work at height commences via a pre-task briefing that confirms conditions on the scaffold or working platform match those described in the assessment. Scaffold inspection records must be maintained on site and signed by the competent inspector. Any scaffold alteration, damage report, or adverse weather event (winds exceeding Beaufort Scale 6) must trigger an immediate inspection and review before work resumes. The MEWP thorough examination certificate must be on site and in date. Any near-miss or fall event must trigger an immediate review of the assessment, suspension of the activity, and notification to the principal contractor before work resumes.}

\example{Five-Step Application}{A bricklayer on a new-build residential development is required to work at 6.5 metres on the second lift of scaffold. Step 1 identifies: fall from scaffold edge; fall through a gap in scaffold boards at the staircase opening; falling brick from the scaffold platform onto operatives below. Step 2 identifies: the bricklayer; the groundworker preparing foundation drainage directly below the scaffold; delivery drivers accessing the hoarding gate beneath the scaffold. Step 3 rates inherent fall risk as Very High (25/25). Step 4 implements: double guardrail, midrail, and toe-board on all open scaffold edges; fully-boarded platform with no gaps; brick guards fitted to the outside of the scaffold lift above the public footpath; the groundworker and delivery access route relocated away from the scaffold footprint; weekly scaffold inspections by a CISRS-certificated inspector. Residual risk is assessed as Low (4/25). Step 5 mandates daily pre-task scaffold check before the bricklayer ascends.}

%---------------------------- SECTION ----------------------------
\section{Applying the hierarchy of controls}

\paragraph{The Working at Height Regulations 2005 embed the hierarchy of controls directly into the regulatory duty structure: avoid, then prevent, then minimise. This means the assessment must demonstrate in writing that avoidance was considered before prevention measures were specified, and that collective prevention was considered before personal protective equipment was selected.}

\begin{itemize}
  \bitem{Elimination}{Avoid work at height entirely: use ground-level prefabrication and crane lift of pre-assembled modules; specify long-reach tools for gutter cleaning and inspection; use drones for roof and facade surveys; design parapets and maintenance access into building designs to facilitate future work at height from a protected position.}
  \bitem{Substitution}{Replace higher-risk access methods with lower-risk alternatives: replace a ladder with a podium step or tower scaffold for tasks of more than thirty minutes duration; replace a tower scaffold with a MEWP where ground conditions permit; replace edge-working on a fragile roof with a purpose-built access system incorporating fall-arrest netting.}
  \bitem{Engineering controls}{Full-boarded scaffold platforms with double guardrail, midrail, and toe-board on all open edges; safety nets, airbags, or soft-landing systems beneath the working platform; edge protection and temporary handrails at all leading edges; covered pedestrian walkways and exclusion zones beneath elevated work areas; brick guards and debris netting on scaffold lifts adjacent to public areas.}
  \bitem{Administrative controls}{Permit-to-work for access onto fragile roofs; scaffold inspection programme with documented records; competency requirements for scaffold erectors (CISRS card); pre-task briefings before each period of work at height; restriction of ladder use to access only (not as a working platform) for tasks exceeding thirty minutes; prohibition of work at height in wind speeds above Beaufort Scale 6.}
  \bitem{PPE}{Safety harness with inertia reel or energy-absorbing lanyard connected to a confirmed, load-tested anchor point; hard hats for all persons beneath or adjacent to working at height activities; toe-boards and brick guards to supplement edge protection in falling-object risk areas; rescue harness and lowering equipment maintained in accessible location with trained rescue team identified.}
\end{itemize}

%---------------------------- SECTION ----------------------------
\section{Cadence of assessment and review triggers}

\paragraph{Working at height is among the highest-risk activities in construction and demands the most frequent assessment review cadence. The dynamic nature of scaffold structures, the deterioration of access equipment, and the potential for rapid change in ground and weather conditions mean that the review cycle for working at height must be built into the daily site management routine.}

\begin{itemize}
  \bitem{Daily pre-task scaffold check}{Before any operative ascends the scaffold each day, the scaffold supervisor or site manager must confirm that the scaffold is in the condition recorded in the last formal inspection, that all guardrails and toe-boards are in place, and that no unauthorised alterations have been made overnight.}
  \bitem{Seven-day scaffold inspection}{The scaffold must be formally inspected by a competent CISRS-certificated person at a maximum of every seven days, and the inspection record completed, signed, and retained on site before the next period of work at height commences.}
  \bitem{After adverse weather}{Any significant weather event --- high winds, heavy rain, snow, or frost --- must trigger a scaffold inspection before work at height recommences; the Working at Height Regulations 2005 and TG20:21 both specify weather-related inspection triggers.}
  \bitem{After any alteration, damage, or near-miss}{Any alteration to the scaffold structure, any report of damage, any near-miss involving the scaffold, or any event that may have affected the integrity of the structure must trigger an immediate inspection by a competent person before the scaffold is used again.}
  \bitem{At each change of working phase}{When the nature of the work at height changes --- from structure to cladding, from cladding to roofing, from roofing to M\&E installation --- the risk assessment must be reviewed to confirm that the existing scaffold configuration and controls remain suitable for the new work activity.}
\end{itemize}

%---------------------------- SECTION ----------------------------
\section{Common failure modes}

\begin{itemize}
  \bitem{Ladders used as work platforms}{Ladders are access equipment, not work platforms. Using a ladder as a work platform for tasks exceeding thirty minutes duration, or for tasks requiring both hands, is a common and frequently fatal failure. The risk assessment must prohibit ladder use as a work platform and specify the appropriate access equipment.}
  \bitem{Fragile roofs not identified or not protected}{Pre-2000 industrial roofs commonly contain asbestos cement sheets, plastic rooflights, and deteriorated metal decking, all of which may be fragile. Failure to conduct a fragile surface survey before any roof access and to implement edge protection and fragile roof covers is a recurring causal factor in falls through roofs.}
  \bitem{Unprotected edges on scaffold}{Removal of guardrails and toe-boards by trades requiring direct access to the building facade --- bricklayers, window installers, cladders --- without replacement with an alternative fall protection measure is a common and predictable failure that must be addressed in the risk assessment and method statement.}
  \bitem{No rescue plan for harness users}{Specifying harness and lanyard as a fall protection measure without a specific, pre-arranged, and rehearsed rescue plan is a regulatory failure under the Working at Height Regulations 2005 and a practical failure that leaves a fallen operative at risk of suspension trauma.}
  \bitem{Scaffold inspection records not maintained}{Scaffold inspections that take place but are not recorded leave no evidence that the duty was discharged, and inspections that are due but not conducted leave the scaffold in an unknown condition. Both failures are common enforcement findings.}
  \bitem{Overreaching from MEWPs}{Operatives overreaching from MEWP baskets to avoid repositioning the machine is a common cause of MEWP instability and fall events. The risk assessment must address the maximum reach distances from the MEWP basket and the consequences of overloading the machine by unsupported loads.}
\end{itemize}

\example{Failure Pattern}{A cladding installation contractor is working on the fourth floor of a steel-framed commercial building using a tube-and-fitting scaffold. To install the final column of cladding panels, the operative removes the inner guardrail on the scaffold platform to allow direct access to the steel frame. The risk assessment specifies ``guardrails to be maintained at all times''; the method statement makes no provision for the guardrail removal that is a predictable requirement of the cladding installation sequence. The operative works for two hours with no inner guardrail and no alternative fall protection. A principal contractor inspection identifies the situation and issues a stop-notice, but the risk assessment failure --- the failure to anticipate and control a foreseeable working method --- has existed for the entire cladding installation period.}

%---------------------------- CHAPTER SUMMARY ----------------------------
\section{Chapter Summary}

\begin{table}[H]
\centering
\begin{tabularx}{\textwidth}{>{\raggedright\arraybackslash}p{3.2cm}X}
\toprule
\textbf{Assessment Element} & \textbf{Operational Requirement} \\
\midrule
Legal/Professional Basis & Working at Height Regulations 2005; PUWER 1998; LOLER 1998. The regulatory hierarchy --- avoid, prevent, minimise --- must be demonstrated in the assessment. \\
Method & Apply the full five-step process and document inherent/residual risk. \\
Controls & Collective protection (guardrail scaffold, nets, airbags) in preference to personal fall protection (harness); specific rescue plan for all harness use; fragile surface identification and protection; scaffold inspection programme. \\
Cadence & Daily pre-task check; seven-day formal scaffold inspection; inspection after adverse weather, alteration, or near-miss; review at each change of working phase. \\
Assurance & Use audit, incident learning, and governance reporting to verify effectiveness. \\
\bottomrule
\end{tabularx}
\end{table}

\begin{itemize}
  \bitem{Key takeaway 1}{Falls from height are the leading cause of construction fatalities; the regulatory hierarchy of avoid, prevent, and minimise must be applied and documented in every working at height risk assessment.}
  \bitem{Key takeaway 2}{Collective protection measures --- guardrail scaffold, safety nets, airbags --- must be used in preference to personal fall protection equipment wherever reasonably practicable; harness and lanyard is a last-resort control, not a primary solution.}
  \bitem{Key takeaway 3}{Every working at height risk assessment that includes harness use must include a specific, pre-arranged, and documented rescue plan identifying the rescue team, the rescue equipment, and the procedure for each scenario.}
  \bitem{Key takeaway 4}{Scaffold must be erected, altered, and dismantled only by CISRS-competent persons and inspected by a competent person before first use, after any alteration or adverse weather, and at a maximum interval of seven days; inspection records must be retained on site.}
\end{itemize}

\barquote{In Chapter~\ref{chap:Excavation and Groundworks Risk Assessment}, we turn from this assessment to the next domain in the Prudentia framework.}
